\documentclass[11pt,reqno]{amsart}

\usepackage[T1]{fontenc}
\usepackage{lmodern}
\usepackage{microtype}
\usepackage[a4paper,margin=32mm]{geometry}
\usepackage{mathtools}
\usepackage{amssymb}
\usepackage{mathrsfs}
\usepackage{booktabs}
\usepackage{array}
\usepackage{float}
\usepackage{enumitem}
\usepackage{xcolor}
\usepackage{tikz}
\usetikzlibrary{arrows.meta,positioning,shapes.geometric}
\usepackage[colorlinks=true,linkcolor=blue!55!black,citecolor=green!40!black,urlcolor=blue!65!black]{hyperref}
\usepackage{aliascnt}

\newtheorem{theorem}{Theorem}[section]
\newaliascnt{proposition}{theorem}
\newtheorem{proposition}[proposition]{Proposition}
\aliascntresetthe{proposition}
\newaliascnt{lemma}{theorem}
\newtheorem{lemma}[lemma]{Lemma}
\aliascntresetthe{lemma}
\newaliascnt{corollary}{theorem}
\newtheorem{corollary}[corollary]{Corollary}
\aliascntresetthe{corollary}
\theoremstyle{definition}
\newaliascnt{definition}{theorem}
\newtheorem{definition}[definition]{Definition}
\aliascntresetthe{definition}
\newaliascnt{remark}{theorem}
\newtheorem{remark}[remark]{Remark}
\aliascntresetthe{remark}

\usepackage[nameinlink,noabbrev]{cleveref}
\crefname{theorem}{theorem}{theorems}
\Crefname{theorem}{Theorem}{Theorems}
\crefname{proposition}{proposition}{propositions}
\Crefname{proposition}{Proposition}{Propositions}
\crefname{lemma}{lemma}{lemmas}
\Crefname{lemma}{Lemma}{Lemmas}
\crefname{corollary}{corollary}{corollaries}
\Crefname{corollary}{Corollary}{Corollaries}
\crefname{definition}{definition}{definitions}
\Crefname{definition}{Definition}{Definitions}
\crefname{remark}{remark}{remarks}
\Crefname{remark}{Remark}{Remarks}

\newcommand{\Z}{\mathbb{Z}}
\newcommand{\Q}{\mathbb{Q}}
\newcommand{\R}{\mathbb{R}}
\newcommand{\abs}[1]{\lvert#1\rvert}
\newcommand{\norm}[1]{\left\lVert#1\right\rVert}
\newcommand{\cont}[1]{\operatorname{cont}\!\left(#1\right)}
\newcommand{\height}[1]{h_{\mathrm{abs}}\!\left(#1\right)}
\newcommand{\dist}[1]{\left\lVert#1\right\rVert}


\newaliascnt{example}{theorem}
\theoremstyle{definition}
\newtheorem{example}[example]{Example}
\aliascntresetthe{example}
\crefname{example}{example}{examples}
\Crefname{example}{Example}{Examples}
\DeclareMathOperator{\tr}{tr}
\DeclareMathOperator{\ord}{ord}
\DeclareMathOperator{\SL}{SL}
\DeclareMathOperator{\GL}{GL}
\DeclareMathOperator{\arcosh}{arcosh}
\newcommand{\eps}{\varepsilon}
\newcommand{\mc}[1]{\mathcal{#1}}
\frenchspacing



\title[Endpoint costs and balanced near-minimizers]{Endpoint costs and balanced near-minimizers\\under Gaussian reflections}
\author{Denys Dutykh}
\address{Mathematics Department, Khalifa University of Science and Technology, PO Box 127788, Abu Dhabi, United Arab Emirates}
\email{denys.dutykh@ku.ac.ae}
\author{Laurent Vuillon}
\address{Univ. Savoie Mont Blanc, CNRS, LAMA, 73000 Chamb\'ery, France}
\email{laurent.vuillon@univ-smb.fr}
\hypersetup{pdftitle={Endpoint costs and balanced near-minimizers under Gaussian reflections},
 pdfauthor={Denys Dutykh and Laurent Vuillon},
 pdfsubject={Balanced finite paths, Euclidean cost and a norm-89 reflection theorem},
 pdfkeywords={Markoff numbers, Gaussian integers, continued fractions, weighted graphs, balanced words, Sturmian words, return words}}
\subjclass[2020]{Primary 68R15; Secondary 05C22, 11D25, 11J70, 68Q45}
\keywords{Markoff numbers, Gaussian integers, continued fractions, weighted graphs, balanced words, Sturmian words, return words}
\begin{document}
\begin{abstract}
A comparison of Euclidean costs along a finite graph of
continued-fraction states depends on where the walk begins and ends as
much as on the cycles it traverses. We prove a quantitative stability
theorem for finite graphs carrying nonnegative integer initial, edge and
terminal weights: if every cyclic component of the zero-weight subgraph
has a single letter frequency, then the heavy-letter count of any
balanced accepted word lies within an explicit multiple of its completed
cost from that finite set of frequencies, the multiple being read off the
condensation graph. Balance is what allows the finite set to replace its
convex hull. The compatible slopes of a finite word are identified with the
projection of a face of the arrangement by which Berstel and Pocchiola
enumerated the factors of Sturmian words, which also measures how thin the
arithmetic input is: of the balanced words of length $n$, only $O(n)$ of the
roughly $n^3/\pi^2$ arise from a Markoff occurrence.

We apply the theorem to the Gaussian reflection attached to $89=8^2+5^2$,
acting on primitive sources whose norm has $89$-adic valuation exactly
one. On an unrestricted paired-digit language this reflection never
decreases subtractive-Euclidean cost, and we determine every word whose
cost it preserves. Consequently every distinct reflected target of a
genuine reduced Markoff source has strictly larger cost, by an amount
that grows linearly with the denominator of the Farey slope outside any
neighborhood of four critical slopes. An explicit infinite family of
genuine sources realizes this growth. The proof of the inequality at $89$
is in part computer-assisted; the stability theorem and the infinite
family are not. These comparisons concern the complete norm-$89$ flip and
do not establish global Markoff uniqueness.
\end{abstract}
\maketitle
\section{Introduction}
\label{sec:introduction}

The Markoff equation\footnote{We use the spelling \emph{Markoff} for the
equation and its numbers, and refer to the uniqueness question as the
Markoff--Frobenius conjecture.}
$$
 a^2+b^2+c^2=3abc
$$
has a tree of positive integer solutions, generated from $(1,1,1)$ by
permutations and by the replacement $(a,b,c)\mapsto(a,b,3ab-c)$. The
Markoff--Frobenius conjecture asserts that the largest entry of a solution
determines the other two. It remains open, and a recurring strategy is to
attach to each occurrence of a largest entry in the tree an invariant finer
than the entry itself, rigid enough to separate two occurrences that share
their numerical value. The invariant studied here is a cost.

Each occurrence carries a reduced fraction $P/Q$, its Farey slope, and
determines a primitive representation of its label as a sum of two squares.
The subtractive Euclidean algorithm assigns an integer cost to such a pair:
the number of subtractions needed to bring the two coordinates to their
common divisor, or equivalently the sum of the partial quotients of their
ratio, minus one. For the occurrence pair that cost is exactly $P+Q-2$.
The natural question is whether the same integer admits a cheaper primitive
representation, and the comparison is subtle: a Gaussian reflection
preserves the norm while changing many partial quotients at once.

Reading the continued fraction of a source and applying the reflection is a
walk through a finite set of integer matrices, along which the cost
difference accumulates. Nonnegative weights on the cycles of such a graph
control repeated motion, but a finite walk also begins and ends somewhere,
and its initial and terminal charges survive any rebalancing of the edge
weights. Our first result concerns exactly that situation; the second
applies it to one Gaussian prime block.

\subsection{The two results}

The general theorem concerns a finite graph reading one binary letter on
each edge. Suppose an exact potential makes all initial, edge and terminal
weights nonnegative integers, and suppose each cyclic component of the
zero-weight subgraph carries a single letter frequency. Denote the finite
set of those frequencies by $\Theta$, assumed nonempty. For a balanced
accepted word of length $n$ with $m$ heavy letters and completed cost
$\Delta$, we prove
\begin{equation}
 \min_{\theta\in\Theta}|m-\theta n|<K(\Delta+1),
 \label{eq:intro-stability}
\end{equation}
where $K$ is obtained explicitly from the condensation graph and from
elementary count potentials inside its cyclic components. No
irreducibility, determinism, or bound on the input length is required.

The role of balance is precise. An unrestricted word can spend a long time
in one recurrent component and then a long time in another, so that its
overall frequency lies strictly between theirs. A balanced word cannot make
such a macroscopic mixture unless the two frequencies are already close.
This is what turns a statement about the convex hull of $\Theta$ into the
finite-set estimate \eqref{eq:intro-stability}.

Balance enters through a single compatible slope, valid simultaneously for
every factor. That slope has a geometric meaning: the pairs consisting of a
slope and an intercept that produce a given finite word form a convex region
of the plane, and the compatible slopes are exactly its shadow on the slope
axis. These regions are the faces of the arrangement through which Berstel
and Pocchiola counted the finite factors of Sturmian words, so the estimate
rests on one coordinate of a two-parameter picture.

Our arithmetic application fixes the Gaussian block $89=8^2+5^2$. If a
primitive source has norm divisible by $89$ exactly once, there is a single
complementary allocation of this pair of conjugate Gaussian primes, modulo
units and conjugation; we call the resulting comparison the complete
norm-$89$ flip. For the entire language $q=\varepsilon\chi(w)2$, where
$\chi$ doubles each digit, we prove that the flip cannot lower the cost,
and we determine every word on which the cost is preserved:
\begin{equation}
 q=111\chi(s)2,\qquad s\in(12111\mid21111)^*.
 \label{eq:intro-equality}
\end{equation}
The target exchanges the two five-letter tokens. The reduced-count
condition satisfied by a genuine Markoff occurrence excludes every nonempty
token sequence, so every distinct genuine target has strictly larger cost.

The same finite graph has critical frequencies
\begin{equation}
 \Theta_{89}=\{0,1/5,1/3,1\}.
 \label{eq:intro-critical}
\end{equation}
For a genuine source of Farey slope $P/Q$, the quantitative consequence is
\begin{equation}
 \min_{\theta\in\Theta_{89}}|P-\theta Q|\leqslant118\Delta+121.
 \label{eq:intro-farey}
\end{equation}
The excess therefore grows linearly with $Q$ on every closed interval of
slopes disjoint from $\Theta_{89}$. We construct an infinite sequence of
genuine sources to which this conclusion applies: their norms have $89$-adic
valuation exactly one, and both Farey coordinates tend to infinity.

\subsection{Background and contribution}

Finite transducers for homographic transformations of continued fractions
originate in Raney's work, and their matrix-state interpretation continues
to be useful in arithmetic applications
\cite{Raney1973,DongDupontODorneyWaitkus2025}. Potential transformations of
weighted automata, including the accompanying changes at initial and
terminal states, are standard \cite{Mohri2009}. Finiteness of an automaton
carries arithmetic information in other settings as well: over a finite
field, the power series that are algebraic are precisely those whose
coefficients a finite automaton can generate \cite{ChristolKamaeMendesFranceRauzy1980}.
We make no use of that equivalence here, and return to it in
\cref{sec:conclusion}. Critical classes and the
eventual concentration of optimal paths are central in min-plus spectral
theory \cite{Quadrat1997}, and the relation between cycle frequencies and
convex rotation sets has a well-developed graph-theoretic interpretation;
see \cite[Section 3]{BochiZhang2016}. We use these mechanisms explicitly and
attribute them where they occur. What \eqref{eq:intro-stability} adds is a
balance condition imposed on individual finite inputs, which replaces the
convex hull of the critical frequencies by the frequencies themselves.

The general statement proved here is the finite balanced-path estimate
\eqref{eq:intro-stability}, with its endpoint hypotheses and an explicit
condensation constant. We also record how narrow the arithmetic input is:
of the roughly $n^3/\pi^2$ balanced words of length $n$, at most $3n+3$ are
cores of genuine occurrences. The norm-$89$ application adds the endpoint
inequality, the complete equality language and the quantitative consequence
for genuine occurrences. We make no priority claim for the underlying
potential method, for the classical balance lemma, or for the
correspondence between central words and Markoff numbers.

The proof of the endpoint inequality at $89$ is in part computer-assisted:
it rests on finitely many integer inequalities attached to a finite graph,
verified in exact arithmetic. \Cref{sec:endpoint} states precisely which
finite assertion is verified and why it implies a statement about words of
arbitrary length. The stability theorem of \cref{sec:balanced-stability} and
the infinite family of \cref{sec:markoff-application} involve no computation
over words or labels.

The result compares a source with its complete norm-$89$ flip. A norm
containing other Gaussian prime blocks can have further primitive
representations, and no comparison with all of them is asserted here; nor
is the critical set \eqref{eq:intro-critical} claimed to be sharp among
genuine Markoff sources. These boundaries leave a concrete question for
further work: which recurrent zero-cost components can a central source
approach while retaining its arithmetic endpoint conditions?

\section{Euclidean cost and a complete Gaussian flip}
\label{sec:cost-reflections}

\subsection{The cost convention}

For a nonzero positive vector $v=(y,x)^{\mathsf T}$, let $\ell(v)$ be the
number of unit subtractions in the algorithm that repeatedly subtracts
the smaller coordinate from the larger and stops when the coordinates
are equal. This definition is unchanged by interchanging coordinates or
multiplying both by a positive integer. For a vector with nonzero signed
coordinates, we apply it to their absolute values.

Write
\begin{equation}
 M(j)=\begin{pmatrix}j&1\\1&0\end{pmatrix},\qquad
 M(q_1\cdots q_r)=M(q_1)\cdots M(q_r),\qquad
 e_1=\begin{pmatrix}1\\0\end{pmatrix}.
 \label{eq:digit-matrix}
\end{equation}
For a positive continued-fraction word $q$ with final digit at least two,
the first column of $M(q)$ is a primitive vector $(y,x)^{\mathsf T}$ with
$y>x>0$. The usual canonical finite expansion of $y/x$ is $[q_1;\ldots,q_r]$.
Division with remainder gives
\begin{equation}
 \ell(M(q)e_1)=q_1+\cdots+q_r-1.
 \label{eq:euclidean-cost}
\end{equation}
Indeed, the successive divisions make $q_i$ subtractions until the last
division, where the algorithm stops at equal coordinates and makes one
fewer. The final minus one is part of our convention throughout.

Let $\chi$ double each digit, so that $\chi(121)=112211$. Our regular source
language is
\begin{equation}
 q=\varepsilon\chi(w)2,\qquad
 \varepsilon\in\{\varnothing,1,2\},\quad w\in\{1,2\}^*.
 \label{eq:regular-source}
\end{equation}
An empty optional digit contributes zero to a digit sum. If $w=j_1\cdots j_n$,
then
\begin{equation}
 \ell(M(q)e_1)=\varepsilon+2\sum_{i=1}^n j_i+1.
 \label{eq:paired-cost}
\end{equation}
The language in \eqref{eq:regular-source} is larger than the genuine
Markoff source language. This enlargement makes the endpoint theorem
stronger, while genuine sources will later supply balance and coprime
Farey counts.

\subsection{The two reflection numerators}

The rational prime $89$ splits in the ring $\Z[i]$ of Gaussian integers.
Fix once and for all the Gaussian prime above it given by
\begin{equation}
 \pi=8+5i,\qquad \bar\pi=8-5i,\qquad \pi\bar\pi=89.
 \label{eq:gaussian-prime}
\end{equation}
The four associates $\pm\pi,\pm i\pi$ generate the same ideal, and
$\bar\pi$ generates the other prime above $89$. Every choice made below is
expressed through this one Gaussian integer, which makes the construction
rigid: the coordinate pair $(8,5)$ of $\pi$ is itself the primitive vector
of the source $q=1112$, identified in \cref{cor:genuine-strict} as the
unique genuine source of vanishing excess. The two integral matrices
\begin{equation}
 R_0=\begin{pmatrix}39&80\\80&-39\end{pmatrix},\qquad
 R_1=\begin{pmatrix}80&39\\39&-80\end{pmatrix}
 \label{eq:reflections}
\end{equation}
satisfy $R_\nu^{\mathsf T}R_\nu=89^2I$ and $\det R_\nu=-89^2$.
Under the identification $(y,x)^{\mathsf T}\leftrightarrow y+ix$, they act
as $z\mapsto\pi^2\bar z$ and $z\mapsto i\bar\pi^2\bar z$, respectively.

For a nonzero integral vector $u$, write $\cont{u}$ for the greatest common
divisor of its coordinates, taken positive. The next lemma explains why
the exact endpoint content, rather than divisibility alone, matters.

\begin{lemma}[Exact primitive endpoint]
\label{lem:primitive-flip}
Let $v=(y,x)^{\mathsf T}$ be primitive, with $x,y>0$, and let
$N=x^2+y^2$. Then $\ord_{89}(N)=1$ if and only if exactly one of $R_0v,R_1v$
has content $89$. For this branch, the vector $R_\nu v/89$ is primitive,
has nonzero coordinates and norm $N$. Modulo signs and coordinate order,
it represents the complete flip of the Gaussian prime pair above $89$.
\end{lemma}

\begin{proof}
The prime $89$ splits as $\pi\bar\pi$ in $\Z[i]$. Since $v$ is primitive,
$z=y+ix$ cannot be divisible by both $\pi$ and $\bar\pi$. If $\pi^e$ divides
$z$ exactly, with $e>0$, then the valuations of $\pi^2\bar z$ at
$(\pi,\bar\pi)$ are $(2,e)$. Its rational content at $89$ is therefore
$89^{\min(2,e)}$. The other branch has content prime to $89$. The roles
reverse when $z$ is divisible by $\bar\pi$. When neither divides $z$,
neither branch has content divisible by $89$.

No rational prime other than $89$ can divide a reflected content: the
matrix is invertible modulo every such prime. Since $N=z\bar z$ and only
one of $\pi,\bar\pi$ divides $z$, the exponent $e$ is exactly
$\ord_{89}(N)$. Thus content exactly $89$ is equivalent to $\ord_{89}(N)=1$. Division gives a primitive norm-$N$ vector.
If one coordinate vanished, $N$ would be a square, contradicting
$\ord_{89}(N)=1$.

For example, if $z=\pi t$, the compatible branch divided by $89$ is
$\pi\bar t$. Its conjugate is $\bar\pi t$, which exchanges precisely the
allocation of $\pi$ and $\bar\pi$. Units and conjugation account for the
irrelevant signs and coordinate order. This is the stated complete flip.
\end{proof}

For the branch selected by \cref{lem:primitive-flip}, define the completed
cost excess
\begin{equation}
 \Delta=\ell(R_\nu v/89)-\ell(v).
 \label{eq:completed-excess}
\end{equation}
The normalization by $89$ does not change the subtractive cost. It does,
however, decide whether the terminal vector is the primitive
representation being compared. If $\ord_{89}(N)\geqslant2$, the compatible
raw content is $89^2$ and the present theorem does not apply.

\section{Genuine Markoff sources and their balanced cores}\label{sec:markoff-sources}

The reflection machine will accept continued fractions from a larger
language than the one supplied by Markoff occurrences. We now identify
the genuine sources inside that language. The distinction matters:
balance controls long factors, whereas the exact occurrence also fixes
the initial phase and the terminal digit.

\subsection{The mechanical word and the Cohn normalization}

Fix relatively prime integers $0<P<Q$. The lower Christoffel word
$C_{P/Q}=c_0\cdots c_{Q-1}$ is defined by
\begin{equation}\label{eq:source-mechanical}
 c_i=b\quad\Longleftrightarrow\quad
 \left\lfloor\frac{(i+1)P}{Q}\right\rfloor
 -\left\lfloor\frac{iP}{Q}\right\rfloor=1.
\end{equation}
It has $P$ letters $b$ and $Q-P$ letters $a$; thus $Q$ is its
total length. Its first letter is $a$ and its last is $b$, so write
$C_{P/Q}=azb$. The interior $z$ is a palindrome. Indeed, for
$1\leqslant j\leqslant Q-1$, coprimality gives
$$
 \left\lfloor\frac{(Q-j)P}{Q}\right\rfloor
 =P-1-\left\lfloor\frac{jP}{Q}\right\rfloor.
$$
Applying this identity to two consecutive interior indices shows that
$c_i=c_{Q-1-i}$ for $1\leqslant i\leqslant Q-2$.
Such interiors are called central words; see
\cite{BertheDeLucaReutenauer2008} for their standard word-theoretic
description.

Put
\begin{equation}\label{eq:source-cohn-matrices}
 A=M(1)^2=\begin{pmatrix}2&1\\1&1\end{pmatrix},\qquad
 B=M(2)^2=\begin{pmatrix}5&2\\2&1\end{pmatrix},
\end{equation}
and let $\mu$ be the word morphism with $\mu(a)=A$ and $\mu(b)=B$.
The Cohn--Christoffel correspondence identifies the normalized
Markoff occurrence at $P/Q$ with the label
\begin{equation}\label{eq:source-label}
 m_{P/Q}=\frac13\operatorname{tr}\mu(C_{P/Q})
        =\mu(C_{P/Q})_{12}.
\end{equation}
This is the classical correspondence with Markoff triples, not an
assertion that distinct fractions have distinct labels
\cite{Reutenauer2009}.
The two initial Markoff values $1$ and $2$ lie outside this proper-word
normalization and will not be needed below.

\subsection{The primitive column and its exact cost}

Use the digit product $M(q)$ from \eqref{eq:digit-matrix}. Identify
$a,b$ with the numerical input letters $1,2$, so that $\chi(a)=11$ and
$\chi(b)=22$. Reversal is denoted by a tilde.

\begin{proposition}[The genuine source dictionary]\label{prop:genuine-source-dictionary}
For every reduced $0<P<Q$, there is a canonical continued-fraction
word
\begin{equation}\label{eq:source-half}
 q=\varepsilon\chi(w)2,\qquad
 \varepsilon\in\{\varnothing,1,2\},
\end{equation}
whose first column $v=(y,x)^{\mathsf T}$ satisfies
$y>x>0$, $\gcd(x,y)=1$, and $x^2+y^2=m_{P/Q}$.
The core $w$ is a finite balanced binary word. Moreover,
\begin{equation}\label{eq:source-counts}
 \begin{aligned}
 P&=|q|_2=1+\mathbf1_{\{\varepsilon=2\}}+2|w|_b,\\
 Q&=|q|+1=2+\mathbf1_{\{\varepsilon\ne\varnothing\}}+2|w|.
 \end{aligned}
\end{equation}
and its subtractive Euclidean cost is
\begin{equation}\label{eq:source-cost}
 \ell(v)=\sum_{j=1}^s q_j-1=P+Q-2.
\end{equation}
The complete word is determined by the central palindrome through
\begin{equation}\label{eq:source-full-palindrome}
 \widetilde q\,q=2\chi(z)2.
\end{equation}
\end{proposition}

\begin{proof}
Split the central palindrome at its middle, in exactly one of the forms
$$
 z=\widetilde w\,w,\qquad
 z=\widetilde w\,a\,w,\qquad
 z=\widetilde w\,b\,w.
$$
The corresponding choices are $\varepsilon=\varnothing,1,2$.
Since $A$ and $B$ are symmetric, in all three cases
\begin{equation}\label{eq:source-half-factorization}
 \mu(z)=H^{\mathsf T}H,\qquad
 H=\begin{cases}
       \mu(w),&\varepsilon=\varnothing,\\
       M(\varepsilon)\mu(w),&\varepsilon\in\{1,2\}.
    \end{cases}
\end{equation}
Let $e=(2,1)^{\mathsf T}=M(2)(1,0)^{\mathsf T}$. The first row of
$A$ is $e^{\mathsf T}$ and the second column of $B$ is $e$.
Consequently
$$
 \mu(azb)_{12}=e^{\mathsf T}\mu(z)e=\|He\|^2.
$$
For completeness, if
$\mu(z)=\begin{pmatrix}a&b\\b&d\end{pmatrix}$, direct multiplication
gives
$$
 \mu(azb)_{12}=4a+4b+d,
 \qquad \operatorname{tr}\mu(azb)=3(4a+4b+d),
$$
which also verifies the equality of the two label conventions in
\eqref{eq:source-label}.

The matrix $H$ is integral and unimodular. Therefore $v=He$ is
primitive. Its coordinates are positive, with the first larger than
the second: every digit matrix $M(j)$ sends a positive column to one
with this ordering. By construction $HM(2)=M(q)$, proving the norm
claim. The last digit of $q$ is $2$, so $q$ is the canonical finite
continued fraction of $y/x$. This also includes the empty center,
where $q=2$ and $v=(2,1)^{\mathsf T}$.

Reversing the digit word in each of the three middle cases gives
\eqref{eq:source-full-palindrome} literally. In particular, the full
word has length $2+2|z|=2Q-2$ and contains
$2+2|z|_b=2P$ digits equal to $2$. Its two reversed halves have equal
length and equal digit counts. Hence $|q|=Q-1$ and $|q|_2=P$,
which gives \eqref{eq:source-counts}.

Equation \eqref{eq:euclidean-cost} gives $\ell(v)=\sum q_j-1$.
Since all digits are $1$ or $2$, their sum is $|q|+|q|_2=Q-1+P$,
proving \eqref{eq:source-cost}.

Finally, a binary word is \emph{balanced} if any two of its contiguous
factors of equal length have numbers of $b$'s differing by at most one.
For a length-$r$ factor of $C_{P/Q}$ starting at $i$, the number of
$b$'s is, by telescoping \eqref{eq:source-mechanical},
$$
 \left\lfloor\frac{(i+r)P}{Q}\right\rfloor
 -\left\lfloor\frac{iP}{Q}\right\rfloor.
$$
For fixed $r$ this is either $\lfloor rP/Q\rfloor$ or
$\lceil rP/Q\rceil$. Thus $C_{P/Q}$ is balanced, as is every
contiguous factor. The right half $w$ of $z$ is one such factor.
\end{proof}

The proof explains both roles of the source grammar. Pairing the digits
produces an integer Gram square and an exact cost; the mechanical word
supplies balance. No cyclic balance of an arbitrary central half is
asserted. In particular, a repeated segment of a finite core must be
analyzed with its actual initial and terminal states.

Conversely, the conditions that $q$ has the form
\eqref{eq:source-half}, that $w$ is balanced, and that the counts in
\eqref{eq:source-counts} are coprime are only necessary conditions in
this paper. An exact recognition test must also compare
$\widetilde q\,q$ with $2\chi(z_{P/Q})2$, where $z_{P/Q}$ is the
interior computed from \eqref{eq:source-mechanical}. If this equality
holds, the first-column construction proves that the source is exactly
the occurrence at $P/Q$. The reflection arguments below apply to a
larger language when this final recognition test is not imposed.

\subsection{How thin the genuine language is}

Balance alone is a weak constraint compared with being a genuine core. The
classical enumeration of the finite factors of Sturmian words makes the
comparison exact. Those factors are precisely the balanced words, and there
are $1+\sum_{i=1}^{n}(n-i+1)\varphi(i)$ of length $n$, with $\varphi$ the
Euler function; Mignosi proved this, and Berstel and Pocchiola gave a
geometric proof by duality and Euler's relation for planar graphs
\cite{Mignosi1991,BerstelPocchiola1993}.

\begin{proposition}[The genuine cores of a given length]
\label{prop:thin-language}
Fix $n\geqslant1$. The pairs $(\varepsilon,w)$ produced by
\cref{prop:genuine-source-dictionary} with $\lvert w\rvert=n$ are in
bijection with the reduced fractions $P/Q$, $0<P<Q$, having
$Q\in\{2n+2,\,2n+3\}$. There are therefore exactly
\begin{equation}
 \varphi(2n+2)+\varphi(2n+3)\leqslant 3n+3
 \label{eq:genuine-core-count}
\end{equation}
of them, whereas the balanced words of length $n$ number
$1+\sum_{i=1}^{n}(n-i+1)\varphi(i)$, which is asymptotic to $n^3/\pi^2$.
\end{proposition}

\begin{proof}
By \eqref{eq:source-counts} one has
$Q=2+\mathbf1_{\{\varepsilon\ne\varnothing\}}+2n$, so $\lvert w\rvert=n$
forces $Q=2n+2$ when $\varepsilon=\varnothing$ and $Q=2n+3$ otherwise;
conversely both values occur, the parity of $Q$ being the parity of
$\lvert z\rvert=Q-2$ that decides the three middle cases in the proof of
\cref{prop:genuine-source-dictionary}. That proof assigns to each reduced
$P/Q$ exactly one pair, so the map is well defined. It is injective because
the pair recovers the central word, $z=\widetilde w\,w$,
$\widetilde w\,a\,w$ or $\widetilde w\,b\,w$ according as $\varepsilon$ is
$\varnothing$, $1$ or $2$, and distinct reduced fractions have distinct
central words. Counting reduced fractions with a prescribed denominator
gives \eqref{eq:genuine-core-count}; the bound follows from
$\varphi(2n+2)\leqslant n+1$, valid because $2n+2$ is even, together with
$\varphi(2n+3)\leqslant2n+2$. Finally $\sum_{i\leqslant x}\varphi(i)\sim3x^2/\pi^2$,
and summing this once more over $i\leqslant n$ gives the stated asymptotic.
\end{proof}

At length $n$ the reflection machine may therefore be fed any of about
$n^3/\pi^2$ balanced words, of which at most $3n+3$ come from a genuine
occurrence: a proportion $O(n^{-2})$. This is the precise sense in which the
arithmetic language is narrower than the combinatorial one, and it locates
where a sharper form of the results below would have to come from, since
balance is blind to the distinction.

A second consequence of genuineness is available at once, and it removes an
auxiliary choice from the argument of \cref{sec:markoff-application}: the
Farey slope of an occurrence may be used directly as a compatible slope for
its core, with no recourse to the existence statement of
\cref{lem:compatible-slope}.

\begin{corollary}[The Farey slope is compatible]
\label{cor:farey-compatible}
Let $w$ be the core of a genuine reduced occurrence of slope $P/Q$. Then
every nonempty factor $v$ of $w$ satisfies
\begin{equation}
 \bigl\lvert\,\lvert v\rvert_b-(P/Q)\lvert v\rvert\,\bigr\rvert<1.
 \label{eq:farey-compatible}
\end{equation}
\end{corollary}

\begin{proof}
The word $w$ is a factor of $C_{P/Q}$, so by the telescoping computation at
the end of the proof of \cref{prop:genuine-source-dictionary} a factor $v$
of length $r$ has $\lvert v\rvert_b$ equal to $\lfloor rP/Q\rfloor$ or
$\lceil rP/Q\rceil$. Moreover $r\leqslant\lvert w\rvert<Q$ and
$\gcd(P,Q)=1$, so $rP/Q$ is not an integer; each of the two values is
therefore at distance strictly less than one from it.
\end{proof}

\section{Balanced finite paths near critical frequencies}
\label{sec:balanced-stability}

This section is independent of Gaussian arithmetic. It isolates the
mechanism that turns a finite graph with nonnegative costs into a
quantitative statement about the words of small cost it accepts.

\subsection{One slope controls every factor}

A finite binary word is \emph{balanced} if any two of its factors of the
same length have heavy-letter counts differing by at most one. We denote
the heavy letter by $b$ in this section. The following classical form of
finite balance is useful because the same slope works for every factor.

\begin{lemma}[A compatible slope]
\label{lem:compatible-slope}
For every nonempty finite balanced word $w$, there is an irrational
$\beta\in(0,1)$ such that every factor $v$ satisfies
\begin{equation}
 \bigl||v|_b-\beta|v|\bigr|<1.
 \label{eq:factor-discrepancy}
\end{equation}
\end{lemma}

\begin{proof}
For two nonempty factors $u,v$, put $p=|u|$, $q=|v|$, $A=|u|_b$ and
$B=|v|_b$. We first prove
\begin{equation}
 |qA-pB|<p+q
 \label{eq:balance-determinant}
\end{equation}
by induction on $p+q$. If $p=q$, balance bounds the left side by $p$.
Otherwise interchange the factors so that $p<q$, and write $v=v_1v_2$
with $|v_1|=p$. Then
$$
 qA-pB=p(A-|v_1|_b)+(q-p)A-p|v_2|_b.
$$
Balance bounds the first term in absolute value by $p$, and the induction
hypothesis bounds the remaining two terms together strictly by $q$.
This proves \eqref{eq:balance-determinant}, which after division by $pq$
is the factorial-balance inequality of the classical theory of Sturmian
words \cite[Proposition 2.1.7]{Lothaire2002}.

For each nonempty factor $v$, take the open interval with endpoints
$(|v|_b-1)/|v|$ and $(|v|_b+1)/|v|$. Inequality
\eqref{eq:balance-determinant} says that every two intervals overlap.
There are finitely many intervals, so their largest lower endpoint is
strictly smaller than their smallest upper endpoint. Every lower endpoint
is below one and every upper endpoint is positive. The common intersection
with $(0,1)$ is therefore nonempty and open. Any irrational $\beta$ in
this intersection proves \eqref{eq:factor-discrepancy}.
\end{proof}

The passage from that inequality to a single compatible slope is also
classical; see \cite[Chapter~2]{Lothaire2002}. We included the proof to
keep the finite-path argument self-contained. The compatible slope need not equal the empirical density
$|w|_b/|w|$. For instance, $00100010000$ is balanced, but its factor $10001$
has discrepancy $12/11$ relative to that density. The common slope in
\cref{lem:compatible-slope} avoids this finite-word endpoint error.

\subsection{The slope interval as a projection}

The interval constructed in the proof of \cref{lem:compatible-slope} is not
an artefact of that proof. It is the shadow, on the slope axis, of a planar
region attached to $w$, and the discrepancy condition is exactly what
survives when the second parameter is eliminated.

Read $b$ as $1$ and $a$ as $0$, and recall that the \emph{mechanical word}
with slope $\beta$ and intercept $\gamma$ has $i$th letter
$\lfloor\beta i+\gamma\rfloor-\lfloor\beta(i-1)+\gamma\rfloor$.

\begin{proposition}[The compatible slopes are a projection]
\label{prop:slope-projection}
Let $w=w_1\cdots w_n$ be a nonempty binary word, put
$k_i=\lvert w_1\cdots w_i\rvert_b$ for $0\leqslant i\leqslant n$, and set
\begin{equation}
 \mc F(w)=\bigl\{(\beta,\gamma)\in\R^2:\;
   k_i\leqslant\beta i+\gamma<k_i+1,\quad 0\leqslant i\leqslant n\bigr\}.
 \label{eq:parameter-region}
\end{equation}
Then $\mc F(w)$ is the set of parameters whose mechanical word of length $n$
is $w$, and its image under $(\beta,\gamma)\mapsto\beta$ is
\begin{equation}
 J(w)=\bigl\{\beta\in\R:\;
   \bigl\lvert\,\lvert v\rvert_b-\beta\lvert v\rvert\,\bigr\rvert<1
   \text{ for every nonempty factor }v\text{ of }w\bigr\},
 \label{eq:slope-interval}
\end{equation}
an open interval with rational endpoints. It is nonempty precisely when $w$
is balanced, and the slope of \cref{lem:compatible-slope} is any irrational
point of $J(w)\cap(0,1)$.
\end{proposition}

\begin{proof}
The constraint at $i=0$ reads $0\leqslant\gamma<1$, so on $\mc F(w)$ one has
$\lfloor\beta i+\gamma\rfloor=k_i$ for every $i$, and the successive
differences of the $k_i$ are the letters of $w$; conversely those
differences determine the $k_i$ from $k_0=0$. This is the first assertion.

Eliminate $\gamma$. The system \eqref{eq:parameter-region} has a solution
$\gamma$ for a given $\beta$ exactly when
$$
 \max_{0\leqslant i\leqslant n}\bigl(k_i-\beta i\bigr)
 <\min_{0\leqslant j\leqslant n}\bigl(k_j+1-\beta j\bigr),
$$
that is, when $(k_i-k_j)-\beta(i-j)<1$ for all $i\neq j$, the case $i=j$
being vacuous. For $i>j$ the factor $v=w_{j+1}\cdots w_i$ has
$\lvert v\rvert_b=k_i-k_j$ and $\lvert v\rvert=i-j$, so the condition is
$\lvert v\rvert_b-\beta\lvert v\rvert<1$; for $i<j$ the same factor, taken
with the roles reversed, gives $\beta\lvert v\rvert-\lvert v\rvert_b<1$.
The two families together are \eqref{eq:slope-interval}. Being a finite
intersection of the open intervals with endpoints
$(\lvert v\rvert_b\mp1)/\lvert v\rvert$, the set $J(w)$ is an open interval
with rational endpoints.

If $\beta\in J(w)$ and $u,v$ are factors of a common length $p$, then
$\bigl\lvert\lvert u\rvert_b-\lvert v\rvert_b\bigr\rvert<2$, hence at most
one by integrality, so $w$ is balanced; the converse is
\cref{lem:compatible-slope}, whose proof exhibits a point of $J(w)$ in
$(0,1)$. An open interval contains an irrational point.
\end{proof}

\begin{remark}[The two-parameter picture]
\label{rem:duality}
The regions $\mc F(w)$, for $w$ ranging over the balanced words of length
$n$, are the faces of an arrangement of lines in the $(\beta,\gamma)$ plane:
each constraint in \eqref{eq:parameter-region} is a half-plane, and the
bounding lines are dual to the lattice points involved. Berstel and
Pocchiola used exactly this duality, together with Euler's relation, to
count the faces and so to recover Mignosi's enumeration of the finite
factors of Sturmian words, of which there are
$1+\sum_{i=1}^{n}(n-i+1)\varphi(i)$ of length $n$, with $\varphi$ the Euler
function \cite{BerstelPocchiola1993,Mignosi1991}. In that reading,
\cref{prop:slope-projection} says that
\cref{lem:compatible-slope} computes the shadow of one face, and that the
intercept is the coordinate being integrated out. The two parameters play
complementary roles: over a fixed irrational slope the fibre of the
projection is cut into $n+1$ cells, one for each factor of that length of a
Sturmian word of that slope, whereas moving the slope changes the face, and
so the word itself.
\end{remark}

\subsection{Endpoint weights and critical components}

Let $G$ be a finite directed graph. Each edge reads exactly one letter in
$\{a,b\}$; parallel edges and nondeterminism are allowed. Initial states
have weights $I(s)$, terminal states have weights $T(s)$, and edges have
weights $c(e)$. Assume all these weights are nonnegative integers.
The completed weight of an accepted path is
\begin{equation}
 \Delta=I(s_0)+\sum_{e\text{ on the path}}c(e)+T(s_n).
 \label{eq:abstract-cost}
\end{equation}

These weights may result from a potential transformation. If original
weights are $w_0,w(e),f$ and a potential is $\lambda$, the relevant transformation is
\begin{equation}
 I(s_0)=w_0-\lambda(s_0),\quad
 c(e)=w(e)+\lambda(s)-\lambda(t),\quad
 T(s_n)=\lambda(s_n)+f(s_n).
 \label{eq:weight-pushing}
\end{equation}
This is the usual weight-pushing transformation \cite[Section 6.3]{Mohri2009}. The terms telescope, preserving the completed cost. Nonnegative cycle
weights alone do not guarantee nonnegative values of both transformed
endpoint weights.

Let $G_0$ retain the edges of weight zero. Its strongly connected components
form an acyclic condensation graph $\mathcal D$. Call a component
\emph{cyclic} if it contains a directed cycle, including a self-loop.
Suppose each cyclic component $C_i$ admits a reduced rational number
$\theta_i=p_i/d_i$, with $d_i>0$, and an integer function $H_i$ on its vertices such that
\begin{equation}
 d_i\mathbf1_{\{e\text{ reads }b\}}-p_i=H_i(t)-H_i(s)
 \label{eq:count-potential}
\end{equation}
for every internal edge $e:s\longrightarrow t$. Put
\begin{equation}
 D_i=\frac{\max H_i-\min H_i}{d_i},\qquad
 \Theta=\{\theta_i:C_i\text{ cyclic}\}.
 \label{eq:critical-data}
\end{equation}
Summation around a cycle gives $0\leqslant\theta_i\leqslant1$.
Summation along an internal path gives
\begin{equation}
 \bigl||v|_b-\theta_i|v|\bigr|\leqslant D_i.
 \label{eq:internal-count-bound}
\end{equation}

Condition \eqref{eq:count-potential} is equivalent to all directed cycles
in the component having the same letter frequency. To see the converse,
assign to a vertex the sum of $d_i\mathbf1_b-p_i$ along a path from a fixed
root. Two choices give the same sum: append a common return path, then
decompose the resulting closed walks into cycles. This also shows how to
check the condition by a finite list of edge identities.

Give each cyclic vertex of $\mathcal D$ weight $D_i+1$, every other vertex
weight zero and every condensation edge weight one. Define
\begin{equation}
 L=\max_{\gamma\text{ a directed path in }\mathcal D}
 \left(\#\text{edges of }\gamma+
       \sum_{C_i\in\gamma\text{ cyclic}}(D_i+1)\right),\qquad K=L+1.
 \label{eq:condensation-budget}
\end{equation}
Paths consisting of a single vertex are included. The number $L$ is finite
and is obtained by the usual longest-path recursion on an acyclic graph.

It is convenient to allow a bounded number of letters to be read outside
the graph. Fix an integer $h\geqslant0$, and suppose that an initial
prefix and a final suffix of total length at most $h$ are consumed before
entering and after leaving $G$, while \eqref{eq:abstract-cost} still
records the completed cost of the whole word. The case in which no letter
is omitted is $h=0$.

\subsection{The stability theorem}

We isolate first the counting step, because it uses only one property of a
real number $\beta$, namely that it be compatible with $w$ in the sense of
\eqref{eq:factor-discrepancy}, and not the particular slope furnished by
\cref{lem:compatible-slope}. The arithmetic application in
\cref{sec:markoff-application} will supply its own compatible slope.

\begin{lemma}[The budget inequality]
\label{lem:budget}
Assume $\Theta\ne\varnothing$. Let $w$ be a nonempty word of length $n$, all
but at most $h$ of whose letters are read by an accepted path, the
remaining ones forming a prefix and a suffix consumed outside the graph.
Let $\beta$ be a real number with
$\bigl||v|_b-\beta|v|\bigr|<1$ for every nonempty factor $v$ of $w$, and put
$\delta=\operatorname{dist}(\beta,\Theta)$. If $\delta\leqslant1$, then
\begin{equation}
 n\delta\leqslant h+\Delta+(\Delta+1)L.
 \label{eq:budget}
\end{equation}
\end{lemma}

\begin{proof}
Let $k$ be the number of positive-weight edges in the path. Integrality and
nonnegative endpoint weights imply $k\leqslant\Delta$.

Split the path into $k$ positive edges and $k+1$ possibly empty zero-edge
segments. A zero segment visits each condensation vertex at most once.
Inside a cyclic component, its internal subpath reads a factor $v$ of the
full word, and the hypothesis on $\beta$ together with
\eqref{eq:internal-count-bound} implies
$$
 \delta|v|\leqslant|\beta-\theta_i||v|<1+D_i.
$$
A noncyclic component is a singleton without an internal edge. Every
remaining letter crosses a condensation edge and contributes at most one
to $\delta$ times length, since $\delta\leqslant1$. Thus each zero segment
contributes at most $L$. The positive edges contribute at most $k$, and
the omitted boundary letters at most $h$. Consequently
$n\delta\leqslant h+k+(k+1)L$, which is \eqref{eq:budget} because
$k\leqslant\Delta$ and $L\geqslant0$.
\end{proof}

\begin{theorem}[Balanced finite-path stability]
\label{thm:balanced-stability}
Assume $\Theta\ne\varnothing$. Let $w$ be a balanced word all but at most
$h$ of whose letters are read by an accepted path, the remaining ones
forming a prefix and a suffix consumed outside the graph. Then
\begin{equation}
 \min_{\theta\in\Theta}\bigl||w|_b-\theta|w|\bigr|
 <K(\Delta+1)+h.
 \label{eq:stability-main}
\end{equation}
Both counts on the left refer to the whole word. In particular, when no
letter is omitted the bound is $K(\Delta+1)$.
\end{theorem}

\begin{proof}
The empty word is immediate. Otherwise put $n=|w|$, $m=|w|_b$, and choose
$\beta$ from \cref{lem:compatible-slope}; it satisfies
\eqref{eq:factor-discrepancy}, and
$\delta=\operatorname{dist}(\beta,\Theta)\leqslant1$ because $\beta$ and
every member of $\Theta$ lie in $[0,1]$. \Cref{lem:budget} applies. Choosing
a member of $\Theta$ nearest to $\beta$ and applying
\eqref{eq:factor-discrepancy} also to the whole word gives
$$
 \min_{\theta\in\Theta}|m-\theta n|
 <1+n\delta\leqslant h+(\Delta+1)(L+1)
 \leqslant h+K(\Delta+1).
$$
\end{proof}

\begin{corollary}[Linear cost away from the critical set]
\label{cor:linear-cost}
Under the hypotheses of \cref{thm:balanced-stability}, if $n>0$ and
$\operatorname{dist}(m/n,\Theta)\geqslant\eta>0$, then
$$
 \Delta>\frac{\eta n-h}{K}-1.
$$
Consequently, for a fixed $h$, every accumulation point of the empirical
frequencies of balanced accepted words of bounded completed cost and
unbounded length lies in $\Theta$.
\end{corollary}

If $G_0$ has no cyclic component, one must not take a minimum over an empty
set. In this case each zero segment has at most $c-1$ edges, where $c$ is
the number of components. The corresponding conclusion is the direct
length bound $n\leqslant h+c(\Delta+1)-1$.

\subsection{Why the finite set is stronger than its convex hull}

Assume again that $\Theta\ne\varnothing$. Without balance, let
$$
 L_0=\max_{\gamma\subset\mathcal D}
 \left(\#\text{edges of }\gamma+\sum_{C_i\in\gamma\text{ cyclic}}D_i\right).
$$
Assign frequency $\theta_i$ to every internal cyclic-component segment,
and any fixed member of $\Theta$ to the remaining letters. The
length-weighted average of these assigned frequencies belongs to
$\operatorname{conv}\Theta$. Telescoping the internal count potentials
therefore gives, for an arbitrary nonempty input,
\begin{equation}
 n\operatorname{dist}(m/n,\operatorname{conv}\Theta)
 \leqslant h+k+(k+1)L_0.
 \label{eq:convex-hull-bound}
\end{equation}
If all critical frequencies coincide, this already controls the count
without balance.

In general the finite set cannot replace its convex hull. Take a zero-cost
$a$ loop at an initial state, a zero-cost $b$ loop at a terminal state,
and a zero-cost $a$ edge between them. Both endpoint weights are zero.
The critical set is $\{0,1\}$, yet the accepted words $a^{r+1}b^r$ have
cost zero and count distance $r$ from that set. For $r\geqslant2$ these
words are unbalanced. The example identifies exactly the obstruction
removed by \cref{lem:compatible-slope}.

\section{The endpoint inequality at norm 89}
\label{sec:endpoint}

\begin{theorem}[The norm-$89$ endpoint inequality]
\label{thm:endpoint-89}
Let $q=\varepsilon\chi(w)2$ be any word in \eqref{eq:regular-source}, put
$v=M(q)e_1$, and assume $\ord_{89}(\|v\|^2)=1$. Then its complete norm-$89$
flip satisfies $\Delta\geqslant0$.
\end{theorem}

The proof replaces the infinitely many admissible words by the finitely
many states of a reduction graph, and then exhibits on that graph an
integer potential making every step of the cost comparison nonnegative.
The finiteness of the graph is not an empirical stopping observation: it
is forced by the determinant.

\subsection{From row subtraction to weighted transitions}

Suppose a nonnegative integral matrix $S$ has two incomparable rows:
neither dominates the other componentwise. After multiplying by $M(j)^2$,
subtract a dominated row from the dominating row until the rows are again
incomparable. These are the elementary row operations
$$
 (a,b;c,d)\longmapsto(a-c,b-d;c,d)
 \quad\hbox{or}\quad
 (a,b;c,d)\longmapsto(a,b;c-a,d-b).
$$
The notation lists the two rows separated by a semicolon. Every such
subtraction is forced for \emph{every} positive continuation vector:
componentwise domination implies domination after multiplication by that
vector. The rows remain nonnegative, and the sum of the entries decreases,
so this reduction terminates. If $k$ unit subtractions are emitted, assign
the transition the weight
\begin{equation}
 w(S,j)=k-2j.
 \label{eq:edge-weight}
\end{equation}
The subtracted $2j$ is exactly the source cost of the paired digit $jj$.

Row exchange has no effect on Euclidean cost. We therefore identify a
matrix with its row exchange and choose the representative with positive
determinant. Every reduced representative takes the form
\begin{equation}
 S=\begin{pmatrix}r+u&b\\r&b+v\end{pmatrix},\qquad
 r,b\geqslant0,\quad u,v\geqslant1.
 \label{eq:crossed-state}
\end{equation}

\begin{lemma}[State invariants]
\label{lem:state-invariants}
Every state reachable from a seed has the crossed form
\eqref{eq:crossed-state}, content one, and determinant
\begin{equation}
 rv+bu+uv=89^2.
 \label{eq:finite-state-bound}
\end{equation}
In particular $r,b,u,v\leqslant89^2$, so only finitely many states occur.
\end{lemma}

\begin{proof}
A seed arises from $R_\nu M(\varepsilon)$ by reading paired letters and
performing row operations. The digit matrices are unimodular and
$\det R_\nu=-89^2$, so a seed has determinant $\pm89^2$; the choice of the
representative with positive determinant fixes the sign, and expanding
\eqref{eq:crossed-state} gives \eqref{eq:finite-state-bound}. Row
subtraction, row exchange and sign change are unimodular, and so are the
digit matrices acting on the right; each therefore preserves both the
determinant up to sign and the greatest common divisor of the entries.
Since a seed has content one, so does every state reachable from it.
Finally \eqref{eq:finite-state-bound} writes $89^2$ as a sum of three
nonnegative products, and $u,v\geqslant1$, so each of $r,b,u,v$ is at most
$89^2$.
\end{proof}

If the input ends in state $S$, its remaining vector is $Se$, where
$e=(2,1)^{\mathsf T}$. The terminal charge is
\begin{equation}
 f(S)=\ell(Se)-1.
 \label{eq:terminal-charge}
\end{equation}
We accept the terminal exactly when $\cont{Se}=89$. Unimodular row
operations preserve vector content, so \cref{lem:primitive-flip} makes
this the exact primitive endpoint test. The final one subtracted in
\eqref{eq:terminal-charge} is the source cost of its last digit $2$.

\subsection{Resolving every initial sign}

The first row of each $R_\nu$ in \eqref{eq:reflections} is positive, but
the second has opposite signs. Starting with $R_\nu M(\varepsilon)$,
read paired letters until the second row has one weak sign componentwise.
At that point its sign is fixed for every positive continuation. Change
its sign if necessary, perform the forced row reduction, and enter the
finite graph.

This prefix procedure has eleven leaves. \Cref{tab:seeds} gives every
leaf state and its seed weight. A seed that consumes the paired prefix
$p$ and emits $k_0$ subtractions has weight
\begin{equation}
 w_0=k_0-\varepsilon-2\sum_{j\in p}j.
 \label{eq:seed-weight}
\end{equation}
The optional digit is again interpreted as zero when absent.

\begin{table}[tb]
\centering
\caption{Complete sign-resolution seeds. A state $(a,b;c,d)$ has the
displayed first and second rows and positive determinant $89^2$.
The prefix consists of paired input letters. These eleven leaves cover
every continuing word after at most two such letters.}
\label{tab:seeds}
\begin{tabular}{ccclr}
\toprule
Branch & $\varepsilon$ & Prefix & State $S_0$ & $w_0$\\
\midrule
$R_0$ & $\varnothing$ & $1$ & $(121,41;37,78)$ & $-1$\\
$R_0$ & $\varnothing$ & $2$ & $(223,10;33,37)$ & $0$\\
$R_0$ & $1$ & $\varnothing$ & $(119,39;41,80)$ & $-1$\\
$R_0$ & $2$ & $\varnothing$ & $(158,39;121,80)$ & $-2$\\
$R_1$ & $\varnothing$ & $1$ & $(195,37;2,41)$ & $0$\\
$R_1$ & $\varnothing$ & $21$ & $(68,33;67,149)$ & $10$\\
$R_1$ & $\varnothing$ & $22$ & $(119,1;52,67)$ & $9$\\
$R_1$ & $1$ & $1$ & $(43,2;17,185)$ & $4$\\
$R_1$ & $1$ & $2$ & $(127,43;120,103)$ & $0$\\
$R_1$ & $2$ & $1$ & $(233,20;35,37)$ & $3$\\
$R_1$ & $2$ & $2$ & $(271,23;68,35)$ & $7$\\
\bottomrule
\end{tabular}
\end{table}

The prefix tree has sixteen nodes, so besides the eleven leaves there are
five words that stop earlier. They are the empty paired prefix with absent
optional digit in each branch, the empty prefix with optional digit $1$ or
$2$ in branch $R_1$, and prefix $2$ with absent optional digit in branch
$R_1$. Appending the final digit $2$ gives raw contents
$1,1,1,1,1$, respectively, so none is accepted. The seed table therefore
covers every primitive source in the theorem, including short words, and
no infinite sign-resolution branch is left over.

\subsection{The potential and the proof}

Starting from the eleven seeds, close the state set under the two
transitions. The next statement describes the resulting graph. It is a
finite assertion about integers, and it is the one place in this section
where we appeal to a computation.

\begin{proposition}[The closed graph and its potential]
\label{prop:closed-graph}
The closure of the eleven seeds of \cref{tab:seeds} under the two
transitions has $11420$ states and $22840$ labeled edges, of which
exactly $116$ states satisfy the primitive terminal condition. Let $\lambda(S)$
denote the least weighted length of a path from a seed to $S$, each seed
contributing its own weight $w_0$. Then
\begin{equation}
 \begin{aligned}
 w_0-\lambda(S_0)&\geqslant0 &&\text{at all eleven seeds},\\
 w(S,j)+\lambda(S)-\lambda(S')&\geqslant0 &&\text{on all $22840$ edges},\\
 \lambda(S)+f(S)&\geqslant0 &&\text{at all $116$ primitive terminals}.
 \end{aligned}
 \label{eq:potential-inequalities}
\end{equation}
\end{proposition}

\begin{proof}
By \cref{lem:state-invariants} the state set is finite, and every quantity
appearing in \eqref{eq:potential-inequalities} is an integer; the assertion
is therefore a finite list of integer inequalities. We have verified it by
computer, in exact integer arithmetic: the state set is closed under the
two transitions, the minimum defining $d$ is attained and computed at
every state, and the three families are then checked on the closed graph.
The verification was carried out twice, independently, and the resulting
table of values of $\lambda$ accompanies this article as supplementary
material.
\end{proof}

The values of $\lambda$ may be negative; it is their transformed combinations in
\eqref{eq:potential-inequalities} that are nonnegative.

\begin{proof}[Proof of \cref{thm:endpoint-89}]
Let $q$ be an accepted word, with seed state $S_0$ and terminal state
$S_f$. Every operation emitted along its path is an actual unit
subtraction on the reflected continuation: the two entries of $Sc$ can
never be equal at an emitted step, since the dominating row minus the
dominated one would then annihilate the positive continuation vector $c$,
making a row of $S$ vanish against $c$ and forcing $\det S=0$, against
\eqref{eq:finite-state-bound}; the unconsumed terminal cost is
\eqref{eq:terminal-charge}, and the consumed source charges are exactly
\eqref{eq:seed-weight} and \eqref{eq:edge-weight}. Consequently
\begin{equation}
 \begin{aligned}
 \Delta
 &=w_0+\sum_e w(e)+f(S_f)\\
 &=\bigl(w_0-\lambda(S_0)\bigr)
   +\sum_{e:S\to S'}\bigl(w(e)+\lambda(S)-\lambda(S')\bigr)
   +\lambda(S_f)+f(S_f),
 \end{aligned}
 \label{eq:exact-endpoint-sum}
\end{equation}
the inserted potentials telescoping along the path. Every term on the
second line is nonnegative by \cref{prop:closed-graph}. No bound on the
length of $q$ entered the argument, so the inequality holds for arbitrary
finite input length.
\end{proof}

\begin{remark}
\label{rem:role-of-the-computation}
It is worth separating the two contributions. The mathematics lies in
\cref{lem:state-invariants} and in the telescoping identity
\eqref{eq:exact-endpoint-sum}: the first makes the state set finite, the
second turns three families of inequalities on that finite set into a
statement about words of unbounded length. The computation contributes
only the closed graph and the integer potential, and both are recomputed
from the seeds rather than read from a stored table.
\end{remark}

\section{The complete equality language}
\label{sec:equality}

The potential proves more than nonnegativity. Equality in
\eqref{eq:exact-endpoint-sum} forces every transformed weight along the
path to vanish, and this reduces an inequality over $11420$ states to a
language on six of them.

\begin{theorem}[Equality and the reflected word]
\label{thm:equality-language}
Under the hypotheses of \cref{thm:endpoint-89}, equality $\Delta=0$
holds exactly when
$$
 q=111\chi(s)2,\qquad s\in(12111\mid21111)^*.
$$
The target is $111\chi(\iota(s))2$, where $\iota$ exchanges the two
five-letter tokens, preserving their order in the token sequence.
\end{theorem}

\subsection{Exhausting the zero-slack paths}

Retain every edge whose transformed weight is zero. In this graph, keep
the states that are both reachable from a zero-slack seed and able to reach
a zero-slack primitive terminal. The complete calculation gives exactly
the states in \cref{tab:equality-states}. The only surviving seed is
branch $R_0$, optional digit $1$, with empty paired prefix and state $s_0$.
The only surviving terminal is $s_1$.

\begin{table}[htbp]
\centering
\caption{The full equality graph. The last column is the terminal charge
$f(S)$ when the state is a primitive terminal; a dash means the state is
not such a terminal. Only $s_0$ is an initial state and only $s_1$ is accepting.}
\label{tab:equality-states}
\begin{tabular}{clrc}
\toprule
State & Matrix $S$ & $\lambda(S)$ & $f(S)$\\
\midrule
$s_0$ & $(119,39;41,80)$ & $-1$ & --\\
$s_1$ & $(115,37;47,84)$ & $-1$ & $1$\\
$s_2$ & $(89,21;0,89)$ & $0$ & --\\
$s_3$ & $(89,0;68,89)$ & $-1$ & --\\
$s_4$ & $(131,21;47,68)$ & $-1$ & --\\
$s_5$ & $(121,37;41,78)$ & $-1$ & --\\
\bottomrule
\end{tabular}
\end{table}

Every edge is shown in \cref{fig:equality-graph}. From $s_0$, the first
letter must be $1$, reaching $s_1$. The two first-return words at $s_1$
are exactly $12111$ and $21111$. Thus the accepted paired words are
$1(12111\mid21111)^*$. The optional initial digit $1$ converts this
to the asserted source language, which proves necessity.

That the returns at $s_1$ are two in number, and not one or three, is what
makes the equality language a free monoid on two generators, and hence what
makes the count relation $Q=5P$ of \cref{cor:genuine-strict} rigid. The
same numerical signature characterizes Sturmian words, whose every factor
has exactly two return words \cite{Vuillon2001,JustinVuillon2000}. The two notions are
distinct: returns to a factor of an infinite word on one side, first
returns to a state of a weighted graph on the other. We assert no
implication between them. Both tokens above nevertheless carry one heavy
letter in five, which is the frequency $\theta=1/5$ that these same six
states contribute to the critical set of \cref{sec:markoff-application}.

\begin{figure}[tb]
\centering
\begin{tikzpicture}[>=Stealth,
 state/.style={circle,draw,minimum size=8mm,inner sep=1pt},
 every edge/.style={draw,->},every node/.style={font=\small}]
\node[state] (s0) at (0,0) {$s_0$};
\node[state,double,double distance=1.2pt] (s1) at (2.1,0) {$s_1$};
\node[state] (s2) at (4.1,1.15) {$s_2$};
\node[state] (s3) at (4.1,-1.15) {$s_3$};
\node[state] (s4) at (6.2,0) {$s_4$};
\node[state] (s5) at (3.1,-2.7) {$s_5$};
\draw[->] (-1,0)--(s0);
\path (s0) edge node[above] {$1:0$} (s1)
 (s1) edge node[above left] {$1:1$} (s2)
 (s1) edge node[below left] {$2:0$} (s3)
 (s2) edge node[above right] {$2:-1$} (s4)
 (s3) edge node[below right] {$1:0$} (s4)
 (s5) edge[bend left=18] node[below left] {$1:0$} (s0);
\draw[->] (s4.east) .. controls (8,0) and (8,-2.7) ..
 node[right,pos=.52] {$1:0$} (s5.east);
\end{tikzpicture}

\caption{All zero-slack accepted paths at norm $89$. An arrow label
$j:w$ gives the paired input letter and its original edge weight;
the transformed weight is zero on every displayed arrow.
The seed is $s_0$ and the double circle marks the only accepting state
$s_1$. Its two first-return words are $12111$ and $21111$.}
\label{fig:equality-graph}
\end{figure}

\subsection{An exact token exchange}

The literal reflected word follows from small matrix identities. Put
\begin{equation}
 Q_0=M(1)^3,\qquad
 K=\begin{pmatrix}47&84\\68&-47\end{pmatrix},\qquad
 H=\begin{pmatrix}89&68\\0&-89\end{pmatrix}.
 \label{eq:token-intertwiners}
\end{equation}
Direct multiplication gives
\begin{equation}
 R_0Q_0=Q_0K,\qquad KM(2)=M(2)H.
 \label{eq:boundary-intertwiners}
\end{equation}
Writing $G_t=M(\chi(t))$, the two return identities are
\begin{equation}
 KG_{12111}=G_{21111}K,\qquad
 KG_{21111}=G_{12111}K.
 \label{eq:token-exchange}
\end{equation}
Induction over the token sequence therefore yields
\begin{equation}
 R_0M(111\chi(s)2)=M(111\chi(\iota(s))2)H.
 \label{eq:full-token-identity}
\end{equation}
The first column of $H$ is $89e_1$. Hence the reflected first column
is $89$ times the primitive positive continuant column of the displayed
target. The source and target have equal digit sums. This proves
sufficiency and the target formula in \cref{thm:equality-language},
including exact primitive acceptance, for every token sequence.

\begin{corollary}[Strict cost increase for genuine sources]
\label{cor:genuine-strict}
Let $v$ be the canonical source of a reduced Markoff occurrence, with
$\ord_{89}(\|v\|^2)=1$. Every distinct complete norm-$89$ target has cost
at least $\ell(v)+1$. The only genuine equality source is
$q=1112$, with $v=(8,5)^{\mathsf T}$, which the flip fixes.
\end{corollary}

\begin{proof}
If an equality word contains $j$ five-letter tokens, its inferred counts
are
$$
 P=1+2j,\qquad Q=5+10j=5P.
$$
A genuine reduced occurrence has $\gcd(P,Q)=1$ by definition, and
\cref{prop:genuine-source-dictionary} identifies the counts of its word
with that $P$ and $Q$. Therefore $j=0$. The base column is $(8,5)^{\mathsf T}$
and \eqref{eq:full-token-identity} fixes it after division by $89$.
All other genuine sources have positive integral excess.
\end{proof}

The proof uses coprime counts only after exhausting the equality
language. It does not use the converse assertion that a balanced paired
word with coprime counts is automatically a genuine occurrence.

\section{Quantitative penalties for genuine Markoff sources}
\label{sec:markoff-application}

The equality graph describes paths with total excess zero. Near-minimizers
can also visit zero-slack components that have a positive cost of entry
or exit. To understand them, we must examine the entire zero-slack graph,
not only the six states surviving the equality test.

\subsection{The full critical set}

We first record the structure of the zero-slack subgraph.

\begin{table}[tb]
\centering
\caption{Every cyclic zero-slack component at norm $89$.
The frequency $\theta_i$ counts paired letters equal to $2$.
The oscillation $D_i$ is that of the rational count potential
$g_i=H_i/d_i$. Every internal edge satisfies
$\mathbf1_{\{j=2\}}-\theta_i=g_i(S')-g_i(S)$.}
\label{tab:critical-components}
\begin{tabular}{rrrr}
\toprule
Number of components & States per component & $\theta_i$ & $D_i$\\
\midrule
1 & 6 & $1/5$ & $1$\\
2 & 3 & $1/3$ & $2/3$\\
2 & 1 & $0$ & $0$\\
2 & 1 & $1$ & $0$\\
\bottomrule
\end{tabular}
\end{table}

\begin{proposition}[The critical components at norm $89$]
\label{prop:critical-components}
The zero-slack subgraph of \eqref{eq:potential-inequalities} has $12684$
edges and $11411$ strongly connected components, of which exactly seven
are cyclic. Their frequencies and oscillations are those listed in
\cref{tab:critical-components}, and every internal edge of each satisfies
the count-potential identity \eqref{eq:count-potential}. The longest path
in the condensation, weighted as in \eqref{eq:condensation-budget}, has
\begin{equation}
 L=58,\qquad K=59.
 \label{eq:axis-constants}
\end{equation}
\end{proposition}

\begin{proof}
All of this concerns the finite graph of \cref{prop:closed-graph}, so it
is again a finite assertion about integers, and we verified it by
computer in exact arithmetic, twice and independently. A path realizing
the maximum traverses $56$ condensation edges together with the
six-state component, whose vertex weight is two; this accounts for the
value of $L$. The states and internal edges of the seven cyclic
components are listed in the supplementary material.
\end{proof}

\begin{theorem}[Norm-$89$ balanced stability]
\label{thm:axis-stability}
Under the hypotheses of \cref{thm:endpoint-89}, suppose the core $w$ is
balanced. With $n=|w|$, $m=|w|_2$ and
$\Theta_{89}=\{0,1/5,1/3,1\}$,
\begin{equation}
 \min_{\theta\in\Theta_{89}}|m-\theta n|
 <59(\Delta+1)+2.
 \label{eq:axis-core-bound}
\end{equation}
If the source is a genuine reduced occurrence of slope $P/Q$, then
\begin{equation}
 \min\{P,\,Q-P,\,|P-Q/5|,\,|P-Q/3|\}\leqslant118\Delta+121.
 \label{eq:axis-farey-bound}
\end{equation}
\end{theorem}

\begin{proof}
\Cref{prop:closed-graph} supplies the nonnegative integer weights
required by \cref{thm:balanced-stability}, and
\cref{prop:critical-components} gives $L=58$ and $K=59$. At most two paired
letters are consumed in the sign-resolution prefix, so the theorem with
$h=2$ gives \eqref{eq:axis-core-bound}.

For a genuine occurrence the auxiliary slope can be dispensed with. By
\cref{cor:farey-compatible} the Farey slope $P/Q$ is itself compatible with
the core, so \cref{lem:budget} may be applied to $\beta=P/Q$. Its hypothesis
$\delta\leqslant1$ holds with room to spare: the complementary intervals of
$\Theta_{89}$ in $[0,1]$ are $[0,1/5]$, $[1/5,1/3]$ and $[1/3,1]$, whose
half-lengths are $1/10$, $1/15$ and $1/3$, so
\begin{equation}
 \delta=\operatorname{dist}(P/Q,\Theta_{89})\leqslant\tfrac13 .
 \label{eq:delta-third}
\end{equation}
The left side of \eqref{eq:axis-farey-bound} is exactly $Q\delta$. Writing
$\mathbf1_{\ne}$ for the indicator of $\varepsilon\ne\varnothing$, the count
dictionary \eqref{eq:source-counts} gives $Q=2n+2+\mathbf1_{\ne}$, whence
$$
 Q\delta=2n\delta+(2+\mathbf1_{\ne})\delta
 \leqslant2\bigl(h+\Delta+(\Delta+1)L\bigr)+3\delta .
$$
With $h=2$, $L=58$ and \eqref{eq:delta-third} the right side is
$2\bigl(2+\Delta+58(\Delta+1)\bigr)+1=2(59\Delta+60)+1=118\Delta+121$.
\end{proof}

The gain over the route through \eqref{eq:axis-core-bound}, which would
give $118\Delta+124$, is small. What matters is that no slope other than
the occurrence's own enters the argument: \cref{cor:farey-compatible}
replaces an existence statement about balanced words by an identity about
the Farey address, and this is the only place where the narrowness recorded
in \cref{prop:thin-language} is actually used.

\begin{corollary}[Slope-dependent linear penalty]
\label{cor:farey-penalty}
If a genuine source in \cref{thm:axis-stability} satisfies
$\operatorname{dist}(P/Q,\Theta_{89})\geqslant\eta>0$, then
$$
 \Delta\geqslant\frac{\eta Q-121}{118}.
$$
Thus any sequence of genuine sources with bounded norm-$89$ excess and
unbounded denominators can accumulate only at the four critical slopes.
\end{corollary}

\begin{proof}
For each $\theta\in\Theta_{89}$ one has
$|P-\theta Q|=Q\,|P/Q-\theta|\geqslant\eta Q$, so the left side of
\eqref{eq:axis-farey-bound} is at least $\eta Q$; rearranging that
inequality gives the bound. The last assertion is the contrapositive:
along a sequence with $\Delta$ bounded and $Q\to\infty$ the quantity
$(\eta Q-121)/118$ is eventually larger than any bound on $\Delta$, so no
positive $\eta$ can persist.
\end{proof}

The constants arise from the full condensation and were not optimized
using the genuine source language. In particular, the last assertion is
a necessary condition on accumulation points, not an existence theorem
at each critical slope.

\subsection{An infinite arithmetic family with growing excess}

\begin{theorem}[A genuine infinite family]
\label{thm:infinite-family}
There is an infinite sequence of genuine reduced Markoff occurrences
$P_j/Q_j$ whose labels have exact $89$-adic valuation one, for which
$P_j,Q_j\longrightarrow\infty$ and
\begin{equation}
 \frac{P_j}{Q_j}\longrightarrow
 \rho=\frac{1+\Phi}{4+5\Phi},\qquad
 \Phi=\frac{1+\sqrt5}{2}\quad\text{the golden ratio}.
 \label{eq:family-limit}
\end{equation}
For all sufficiently large $j$, their complete norm-$89$ excesses satisfy
\begin{equation}
 \Delta_j\geqslant\frac{Q_j}{1180(4+5\Phi)}-\frac{121}{118}.
 \label{eq:family-penalty}
\end{equation}
\end{theorem}

\begin{proof}
Start with the Markoff triple $(1,34,89)$ and iterate
\begin{equation}
 F(a,b,c)=(b,c,3bc-a).
 \label{eq:vieta-forward}
\end{equation}
The identity $a^2+b^2+c^2=3abc$ is preserved. If $1\leqslant a<b<c$,
then $3bc-a>c$, so the positive triples continue forward in the Markoff
tree and their maxima increase strictly.

Modulo $89^2$, the same map is a permutation because it has the polynomial
inverse
$$
 F^{-1}(a,b,c)=(3ab-c,a,b).
$$
The orbit of the initial residue triple is therefore periodic from its
first term. Let $T>0$ be a period. For every $j\geqslant0$, the last
coordinate of $F^{jT}(1,34,89)$ is congruent to $89$ modulo $89^2$.
It consequently has valuation exactly one at $89$.

The initial labels $34$ and $89$ occur at the adjacent Farey vertices
$1/4$ and $1/5$; their other common neighbour $0/1$ has label one. The
Vieta step is exactly the Farey mediant: the Christoffel word of the
mediant of two Farey neighbours is the concatenation of their Christoffel
words, so $\mu$ of the mediant is the product of the two matrices, and
the Fricke identity $\tr(XY)+\tr(XY^{-1})=\tr X\,\tr Y$ turns this into
$$
 m_{\text{mediant}}=3m_\alpha m_\beta-m_{\text{other}}
$$
for the three labels concerned. Every step of \eqref{eq:vieta-forward}
therefore replaces the older vertex by the mediant of the two retained
vertices. Hence the two Farey columns follow
Fibonacci addition from $(1,4)$ and $(1,5)$. Their determinant has absolute
value one, so all resulting fractions are reduced. Both coordinates grow
without bound, and their ratios tend to \eqref{eq:family-limit}. The
subsequence at multiples of $T$ has the same limit.

This limit lies strictly between $1/5$ and $1/4$. Its nearest critical
slope is $1/5$, and
$$
 \rho-\frac15=\frac{1}{5(4+5\Phi)}>0.
$$
For sufficiently large $j$, the distance of $P_j/Q_j$ from $\Theta_{89}$
is at least half this positive number. Apply \cref{cor:farey-penalty}
with $\eta=1/[10(4+5\Phi)]$ to obtain
\eqref{eq:family-penalty}.
\end{proof}

The argument needs neither the value of $T$ nor a conjecture about prime
values of recurrence sequences. It is explicit in the sense that the
family is described by a rule and not merely shown to exist; the period
$T$ is large, so the members beyond the first are not small numbers. Its modular periodicity selects the
required valuation, while the real Farey dynamics places the source
slopes away from the critical set. These two elementary mechanisms
make the quantitative theorem nonvacuous on genuine sources with
unbounded numerators.

\subsection{Why the equality slope alone is insufficient}

The six-state equality graph has frequency $1/5$, but the full critical
graph also contains the state
$$
 S=\begin{pmatrix}101&19\\19&82\end{pmatrix}
  =\begin{pmatrix}10&1\\1&9\end{pmatrix}^{\!2}.
$$
This state commutes with $M(1)^2$. Exactly two forced row subtractions
return it after one light paired letter, so the loop has zero cost.
A fixed entry and primitive exit give the explicit family
\begin{equation}
 q_j=\chi(21221122112\,1^j\,2122121)2,\qquad j\geqslant0.
 \label{eq:zero-light-family}
\end{equation}
The prefix enters $S$ with cumulative weight $15$. The suffix reaches
$(109,49;43,92)$ with additional weight $20$; its terminal vector is
$89(3,2)^{\mathsf T}$ and its terminal charge is $1$.
Thus every member has exact norm valuation one at $89$ and completed
excess $36$. Its inferred counts are $P=21$ and $Q=38+2j$.
This rules out a cost bound proportional only to $|Q-5P|$ on the full
regular language.

Every displayed core contains both $11$ and $22$, so it is unbalanced.
The family is consequently not a family of genuine Markoff sources and
does not contradict \cref{thm:axis-stability}. It explains why the
zero-slack graph outside the exact equality language cannot simply be
discarded when studying near-minimizers.

\section{Further questions}
\label{sec:conclusion}

The completed norm-$89$ comparison illustrates two distinct sources of
rigidity. The endpoint potential makes cost nonnegative on a large regular
language. Its equality graph then forces a count relation incompatible
with a reduced genuine occurrence, except at the fixed base source.
Finite balance adds a quantitative constraint: long paths with small
completed excess must have frequency close to a recurrent zero-cost
component.

Three questions remain particularly concrete. First, one can seek a
description of endpoint potentials that persists as the Gaussian block
varies. Any such statement must retain both primitive content and the
actual source and terminal charges. The present proof exhibits one explicit potential at $89$ and supplies no
potential uniform in the block.

Second, the graph frequencies $0,1/5,1/3,1$ need not all be attainable as
limits of genuine sources with bounded excess. The exact mechanical
recognition condition is stronger than finite balance and coprime
counts. Incorporating it into the critical-component analysis could
remove some exceptional slopes or improve the constants. The current
theorem makes neither assertion.

Third, minimization on an entire Gaussian fiber requires comparing
allocations involving other prime blocks. Strict increase under one
complete flip is a useful local statement about this arithmetic orbit,
but it does not compare every primitive representation of the same norm.
The infinite family in \cref{thm:infinite-family} provides genuine
instances where this one-block penalty becomes large; understanding
how several blocks interact remains a separate problem.

Two further directions were suggested to us, and both concern the finite
graph rather than the arithmetic. The first is the decomposition used in
the proof of \cref{lem:budget}, which cuts an accepted path at its
positive-weight edges and treats each zero-weight segment separately. That
is a first-return decomposition relative to the zero-slack subgraph, and the
theory of return words gives a sharper vocabulary for such decompositions
than the condensation budget \eqref{eq:condensation-budget} does: the
budget charges every component of a longest condensation path, whereas a
return-word induction would charge only the components a given word
actually revisits. Whether that replaces the constant $K=59$ by something
close to the six-state component alone is the concrete question.

The second is the algebraic reading of finiteness recalled in
\cref{sec:introduction}. Over a finite field, a power series is algebraic
exactly when a finite automaton generates its coefficients
\cite{ChristolKamaeMendesFranceRauzy1980}, and the reflection graph of
\cref{sec:endpoint} is a finite automaton attached to an arithmetic
operation. The analogy stops there, since that equivalence is a
positive-characteristic phenomenon and our labels are ordinary integers; we
do not know whether some series formed from the costs $\Delta$ along the
graph satisfies an algebraic relation, and we have not found a formulation
in which the question is more than suggestive.

The general theorem in \cref{thm:balanced-stability} can be used outside
the Markoff setting. It takes as input a finite graph whose initial, edge
and terminal weights are nonnegative integers and whose zero-weight
cyclic components each carry a single letter frequency, and it returns an
explicit bound on the letter count of every balanced word the graph
accepts. Separating that mechanism from the arithmetic seems to us the
most transferable part of this article.

\raggedbottom
\section*{Supplementary material}
The supplementary material accompanying this article contains the exact
computations used in the proofs of
\cref{thm:endpoint-89,thm:equality-language,thm:axis-stability}: the closed
state graph, its table of potentials, the cyclic components of the zero-slack
subgraph, and instructions for repeating the calculations. The reconstruction
of the graph and of its potential was carried out twice, by implementations
sharing no code. The material also carries independent confirmations of two
results that are proved analytically above and take no computational input:
the letter counts, the cost and the balance of the dictionary of
\cref{prop:genuine-source-dictionary}, together with a direct evaluation of the
excess from its definition, and the polynomial identities behind
\cref{thm:infinite-family}.

\section*{Code and data availability}
The source code and the data supporting this article, with the instructions
that repeat every calculation, are contained in the supplementary material of
this article, which is available, together with that of the other articles of
this series:

{\centering
\url{https://github.com/dutykh/fibonacci-pell-nearest-gaps/}
\par}

\bibliographystyle{amsplain}
\bibliography{references}
\end{document}
